\documentclass[11pt]{article}

\usepackage[utf8]{inputenc}
\usepackage[T1]{fontenc}
\usepackage[margin=1in]{geometry}
\usepackage{amsmath,amssymb,amsthm}
\usepackage{booktabs}
\usepackage{xcolor}
\usepackage{listings}
\usepackage{hyperref}
\usepackage{enumitem}

\hypersetup{colorlinks=true,linkcolor=blue,citecolor=blue,urlcolor=blue}

\newtheorem{theorem}{Theorem}[section]
\newtheorem{lemma}[theorem]{Lemma}
\newtheorem{corollary}[theorem]{Corollary}
\newtheorem{definition}[theorem]{Definition}
\newtheorem{remark}[theorem]{Remark}

\lstdefinelanguage{Lean4}{
  morekeywords={theorem,lemma,def,noncomputable,structure,namespace,end,import,open,
    section,variable,instance,class,inductive,where,deriving,Repr,DecidableEq,
    by,do,let,have,show,from,fun,match,with,if,then,else,calc,exact,intro,
    rfl,sorry,axiom,private},
  sensitive=true,
  morecomment=[l]{--},
  morecomment=[s]{/-}{-/},
  morestring=[b]",
}
\lstset{
  basicstyle=\linespread{0.9}\ttfamily\scriptsize,
  language=Lean4,
  breaklines=true,
  breakatwhitespace=true,
  columns=fullflexible,
  showstringspaces=false,
  keywordstyle=\color{blue!60!black}\bfseries,
  commentstyle=\color{green!40!black},
  stringstyle=\color{orange!70!black},
  frame=single,
  numbers=left,
  numberstyle=\tiny\color{gray},
  tabsize=2,
  extendedchars=true,
  literate=
    {§}{{\S}}1
    {¬}{{$\lnot$}}1
    {²}{{$^{2}$}}1
    {³}{{$^{3}$}}1
    {·}{{$\cdot$}}1
    {¹}{{$^{1}$}}1
    {é}{{\'e}}1
    {–}{{--}}1
    {—}{{---}}1
    {′}{{$'$}}1
    {″}{{$''$}}1
    {⁰}{{$^{0}$}}1
    {⁴}{{$^{4}$}}1
    {₀}{{$_{0}$}}1
    {₁}{{$_{1}$}}1
    {₂}{{$_{2}$}}1
    {₃}{{$_{3}$}}1
    {₇}{{$_{7}$}}1
    {₈}{{$_{8}$}}1
    {ℕ}{{$\mathbb{N}$}}1
    {ℚ}{{$\mathbb{Q}$}}1
    {ℝ}{{$\mathbb{R}$}}1
    {ℤ}{{$\mathbb{Z}$}}1
    {←}{{$\leftarrow$}}1
    {→}{{$\rightarrow$}}1
    {∀}{{$\forall$}}1
    {∃}{{$\exists$}}1
    {∈}{{$\in$}}1
    {∑}{{$\sum$}}1
    {−}{{$-$}}1
    {∘}{{$\circ$}}1
    {≠}{{$\neq$}}1
    {≡}{{$\equiv$}}1
    {≤}{{$\leq$}}1
    {≥}{{$\geq$}}1
    {⌈}{{$\lceil$}}1
    {⌉}{{$\rceil$}}1
    {═}{{=}}1
    {║}{{|}}1
    {╔}{{+}}1
    {╗}{{+}}1
    {╚}{{+}}1
    {╝}{{+}}1
    {⟨}{{$\langle$}}1
    {⟩}{{$\rangle$}}1
    {Σ}{{$\Sigma$}}1
    {α}{{$\alpha$}}1
    {θ}{{$\theta$}}1
    {λ}{{$\lambda$}}1
    {π}{{$\pi$}}1
    {ρ}{{$\rho$}}1
    {ᵢ}{{$_{i}$}}1,
}

\title{Formal Verification of Structural Properties\\
of Cooper's Sporadic Ap\'ery-Like Sequences:\\
A Lean 4 Development}

\author{Xavier Callens \\ \small{Dual-Scale Topological Universe Model Project}}
\date{\today}

\begin{document}
\maketitle

\begin{abstract}
We report a Lean~4 / mathlib formalization of structural properties of three
sporadic Ap\'ery-like integer sequences due to Cooper (2012)~\cite{Cooper2012},
denoted $s_7$, $s_{10}$, $s_{18}$ in the literature's standard labeling
(Gorodetsky 2023~\cite{Gorodetsky2023}). The development covers: (i) sourced,
citation-carrying encodings of the sequences and their defining three-term
recurrences; (ii) their associated order-3 Picard--Fuchs differential operators
in $\theta$-form ($\theta = z\,d/dz$), with machine-checked operator order;
(iii) a symmetric-square API validated against two first-principles examples;
and (iv), the paper's principal new result, a single fully generic Lean proof
that the Cooper operator family satisfies the Almkvist--van Straten
self-adjointness criterion $W \equiv 0$~\cite{AlmkvistvanStraten2021} for
\emph{every} choice of the four defining parameters $(a,b,c,d) \in \mathbb{Q}^4$,
established as a cleared-denominator polynomial identity discharged by the
\texttt{ring} decision procedure and requiring no axioms. We also record a
correction: two previously conjectured polynomial growth bounds for $s_7$ and
$s_{10}$ are false, and we give machine-checked exponential lower bounds
disproving them. All results are stated with an explicit account of what
remains open (in particular, the $\forall n$ recurrence-satisfaction claims for
$s_7$ and $s_{10}$, which are mathematically established in the literature via
Wilf--Zeilberger certificates but not yet formalized). The development builds
with zero \texttt{sorry} and zero non-quarantined axioms against a pinned
mathlib commit; source and build instructions are given in
Section~\ref{sec:repro}, and the complete Lean source discussed here is
reproduced verbatim in Appendix~\ref{app:source}.
\end{abstract}

\tableofcontents

% ════════════════════════════════════════════════════════════════════
\section{Introduction}

Sporadic Ap\'ery-like sequences are integer sequences $u(n)$ satisfying a
three-term holonomic recurrence with polynomial coefficients in $n$, whose
generating functions are periods of families of algebraic varieties and which
exhibit unexpectedly strong integrality and congruence properties, in the
spirit of Ap\'ery's celebrated proof of the irrationality of $\zeta(3)$. The
complete known list comprises 15 such sequences: six due to
Zagier~\cite{Zagier2009}, six due to Almkvist and
Zudilin~\cite{AlmkvistvanStraten2021}, and three due to
Cooper~\cite{Cooper2012}. This paper is concerned with formally verifying
structural properties of Cooper's three sequences, which we denote $s_7$,
$s_{10}$, $s_{18}$ following the labeling of Gorodetsky's unifying
survey~\cite{Gorodetsky2023}.

\paragraph{Motivation.} Beyond their intrinsic number-theoretic interest, these
sequences' associated Picard--Fuchs operators are conjecturally related, via a
classical \emph{symmetric-square} construction, to periods of K3 surfaces. This
relation is of interest to a broader research program (outside the scope of
this paper) investigating physical interpretations of such geometric
structures; the present formalization exists to place the purely mathematical
substrate of that program on a machine-checked footing, independent of and
prior to any physical interpretation. We follow the discipline, standard in
formal mathematics, of stating exactly what is proved and flagging everything
else as open or conjectural.

\paragraph{Contributions.} Concretely, this paper documents a Lean~4 development
that:
\begin{enumerate}[label=(\roman*)]
  \item encodes $s_7$, $s_{10}$, $s_{18}$ and their defining recurrence
    parameters with citations to primary literature (Section~\ref{sec:sequences});
  \item encodes the associated order-3 $\theta$-form Picard--Fuchs operators and
    proves their order is exactly 3 for every parameter choice
    (Section~\ref{sec:theta});
  \item provides a symmetric-square operator API, validated against two
    examples worked from first principles (Section~\ref{sec:symsquare});
  \item proves, as a single fully generic theorem, that the Almkvist--van
    Straten self-adjointness criterion holds identically across the entire
    Cooper operator family --- the paper's main new result
    (Section~\ref{sec:w0});
  \item corrects a previously asserted (and false) polynomial growth bound for
    $s_7$ and $s_{10}$, replacing it with a machine-checked disproof
    (Section~\ref{sec:growth});
  \item documents integrality and the precise, honest boundary of what is
    verified about recurrence satisfaction (Section~\ref{sec:integrality}).
\end{enumerate}

\paragraph{Scope.} This paper reports \emph{only} the arithmetic/algebraic
formalization described above. The host repository additionally contains
exploratory Lean code relating these operators to speculative physical
models (moduli stabilization, chameleon screening); that material rests on
stipulated, non-derived numerical parameters and in one case an axiom that is
self-documented as vacuous, and is explicitly excluded from this report. See
Section~\ref{sec:scope} for a precise statement of this boundary.

% ════════════════════════════════════════════════════════════════════
\section{Background}

\subsection{The Cooper recurrence template}
\label{sec:background-template}

Cooper's three sporadic sequences, together with Zagier's and
Almkvist--Zudilin's, are special cases of a single recurrence template
parameterized by four integers $(a,b,c,d)$:
\begin{equation}
  \label{eq:cooper-recurrence}
  (n+1)^3\, u(n+1) \;=\; (2n+1)\bigl(a n^2 + a n + b\bigr)\, u(n)
    \;-\; n\bigl(c n^2 + d\bigr)\, u(n-1), \qquad u(-1)=0,\ u(0)=1.
\end{equation}
Cooper's three instances are (Table~\ref{tab:params}, sourced from
Gorodetsky~\cite{Gorodetsky2023}, cross-checked against Cooper's original
table~\cite{Cooper2012}):

\begin{table}[h]
\centering
\begin{tabular}{lcccc l}
\toprule
Sequence & $a$ & $b$ & $c$ & $d$ & Closed form \\
\midrule
$s_7$ & 13 & 4 & $-27$ & 3 & $\sum_{k=\lceil n/2\rceil}^{n} \binom{n}{k}^2\binom{n+k}{k}\binom{2k}{n}$ \\
$s_{10}$ & 6 & 2 & $-64$ & 4 & $\sum_{k=0}^{n} \binom{n}{k}^4$ \\
$s_{18}$ & 14 & 6 & 192 & $-12$ & (no verified closed form in this repository; see \S\ref{sec:integrality}) \\
\bottomrule
\end{tabular}
\caption{Cooper's three sporadic sequences and their recurrence parameters.}
\label{tab:params}
\end{table}

\subsection{$\theta$-form Picard--Fuchs operators}
\label{sec:background-theta}

Writing $\theta = z\,\tfrac{d}{dz}$ for the Euler (logarithmic) derivative,
the generic Cooper family corresponds to the order-3 differential operator
\begin{equation}
  \label{eq:cooper-theta}
  L \;=\; \theta^3 \;-\; z(2\theta+1)\bigl(a\theta^2 + a\theta + b\bigr)
    \;+\; z^2\bigl(c(\theta+1)^3 + d(\theta+1)\bigr),
\end{equation}
whose formal power series solution has Taylor coefficients satisfying
\eqref{eq:cooper-recurrence} (Gorodetsky~\cite{Gorodetsky2023}
eq.~(1.7)). The $\theta$-representation is convenient because the operator's
\emph{order} --- the datum this paper certifies --- is simply its
$\theta$-degree, with no need to normalize the leading coefficient to~1.

\subsection{The symmetric-square criterion}
\label{sec:background-c3}

A classical construction associates to a monic order-2 linear differential
operator $L_2 = D^2 + p D + q$ (with $D = d/dz$) an order-3 operator
$\mathrm{Sym}^2(L_2)$ annihilating pairwise products of solutions of $L_2$:
\begin{equation}
  \label{eq:symsquare-formula}
  \mathrm{Sym}^2(L_2) \;=\; D^3 + 3p\,D^2 + \bigl(2p^2+p'+4q\bigr) D + \bigl(4pq+2q'\bigr).
\end{equation}
An order-3 operator $L_3$ that arises this way is called a symmetric square. A
classical necessary and sufficient condition (Almkvist--van
Straten~\cite{AlmkvistvanStraten2021}) for a monic order-3 operator
$D^3 + a_2 D^2 + a_1 D + a_0$ to be \emph{some} symmetric square is the
vanishing of
\begin{equation}
  \label{eq:W-criterion}
  W \;=\; \tfrac13 a_2'' \;+\; \tfrac23 a_2 a_2' \;+\; \tfrac{4}{27} a_2^3
    \;+\; 2a_0 \;-\; \tfrac23 a_1 a_2 \;-\; a_1'.
\end{equation}
Section~\ref{sec:w0} formalizes and proves this criterion for the entire
Cooper family.

% ════════════════════════════════════════════════════════════════════
\section{Formalization methodology}
\label{sec:methodology}

The development is written in Lean~4~\cite{lean4} against a pinned mathlib
commit (Section~\ref{sec:repro}). Three engineering disciplines, enforced by
repository tooling, are relevant to how this paper's claims should be read:

\begin{description}
  \item[No silent \texttt{sorry}.] Incomplete proofs (\texttt{sorry}) are
    permitted only on development branches, never in the material reported
    here; any genuinely open statement is instead recorded as a named,
    explicitly-labeled open goal (Section~\ref{sec:integrality}) rather than
    left as an unproven claim in running code.
  \item[Axiom quarantine.] New \texttt{axiom} declarations are permitted only
    in a dedicated \texttt{Axioms/} directory, each with a docstring
    justification and an entry in a repository-level axiom registry. None of
    the results reported in this paper depend on any axiom; the one axiom
    presently in the registry (Section~\ref{sec:scope}) is unrelated to the
    sequence/operator formalization and is excluded from this report's scope.
  \item[Citation discipline.] Every definition encoding a literature object
    (a recurrence, a coefficient set, an operator template) carries a source
    docstring; numeric parameters are not accepted from model memory but must
    trace to a specific table in a specific cited source.
\end{description}

We also adopt, as a matter of intellectual honesty, an explicit three-tier
epistemic labeling for any claim outside pure Lean-checked mathematics:
\emph{proved} (machine-checked, stated as fact), \emph{checked-but-unproved}
(verified computationally to a stated finite order, not a universally
quantified proof), and \emph{conjectural} (asserted in the literature or by
analogy, not independently verified here). Every result below is labeled.

% ════════════════════════════════════════════════════════════════════
\section{The sequences: encoding and citations}
\label{sec:sequences}

\emph{Status: proved (definitional).}

Table~\ref{tab:params}'s parameters are encoded as a structure
\lstinline{CooperRecurrenceParams} with fields $a,b,c,d : \mathbb{Z}$, and
$s_7$, $s_{10}$ are given \emph{closed-form} definitions as
\texttt{Finset.sum}s of products of binomial coefficients (Listing in
Appendix~\ref{app:cooper-recurrences}), matching the closed forms of
Table~\ref{tab:params}. Both closed forms are independently, exactly
integer-arithmetic verified against 20 golden reference values
(\lstinline{s7_golden_test}, \lstinline{s10_golden_test} in
\texttt{Tests/CooperSequences.lean}, discharged by \texttt{native\_decide} ---
see the disclosure convention below) and against \eqref{eq:cooper-recurrence}
at the first several indices by direct kernel computation. $s_{18}$ has no
verified closed form in this repository (a direct transcription of the
literature's stated closed form disagreed with the recurrence at $n=3$, an
integer-truncation/signed-binomial edge case); it is instead represented by
its sourced recurrence parameters and golden values computed directly from
\eqref{eq:cooper-recurrence}.

\begin{remark}[\texttt{native\_decide} disclosure]
Some numeric checks above involve binomial coefficients too large for the
kernel's default reduction (\texttt{decide}) to evaluate in reasonable time
(e.g.\ $\binom{38}{19}$) and instead use \texttt{native\_decide}, which
delegates evaluation to compiled code. This extends the trusted computing
base to include the Lean compiler, in addition to the kernel; every theorem
using it is tagged accordingly in the source, and we repeat that disclosure
here rather than silently citing such results as pure kernel checks.
\end{remark}

% ════════════════════════════════════════════════════════════════════
\section{$\theta$-form operators and order classification}
\label{sec:theta}

\emph{Status: proved.}

The generic Cooper $\theta$-operator~\eqref{eq:cooper-theta} is encoded as an
element of \lstinline{Polynomial (Polynomial ℚ)} (outer variable $\theta$,
inner variable $z$), parameterized by \lstinline{CooperRecurrenceParams}. This
representation is chosen specifically to avoid monic normalization: dividing
the operator by its leading $\theta$-coefficient would in general introduce
genuine rational-function (as opposed to polynomial) coefficients, since
\texttt{mathlib}'s \lstinline{RatFunc} type does not, at the pinned commit,
carry a differentiation API --- a concrete obstruction encountered and
documented during this development (Section~\ref{sec:w0} explains how the
paper's main result sidesteps it entirely).

We prove:
\begin{theorem}[Operator degree, kernel fact]
\label{thm:degree}
For every \lstinline{p : CooperRecurrenceParams}, the $\theta$-degree of the
encoded operator \lstinline{cooperThetaOperator p} equals exactly 3.
\end{theorem}
This is Lean's \lstinline{cooperThetaOperator_natDegree}; the proof expands
the operator into explicit collected $\theta$-coefficients
(\lstinline{cooperThetaOperator_eq}, a \texttt{ring} identity) and applies
\lstinline{compute_degree!}, discharging the side condition that the leading
coefficient $1-2az+cz^2$ is nonzero (its constant term is~1, for every
parameter choice). The analogous order-2 statement for the Zagier template is
proved by the same method. Theorem~\ref{thm:degree} replaces two axioms that
were present in an earlier version of this repository
(\lstinline{empirical_s7_degree}, \lstinline{empirical_S12_degree}), which
asserted only the vacuous statement ``there exists a cubic polynomial of
degree 3'' --- true of \emph{any} cubic and hence carrying no information
about the specific encoded operator. Theorem~\ref{thm:degree} is the
non-vacuous replacement: it is a fact about the concrete, sourced coefficient
data.

We stress what Theorem~\ref{thm:degree} does \emph{not} claim: minimality of
this operator for its sequence (i.e.\ that no lower-order operator annihilates
the same generating function) and any identification of the operator's
periods with an actual geometric object are separate, unformalized claims not
addressed by this paper.

% ════════════════════════════════════════════════════════════════════
\section{The symmetric-square API}
\label{sec:symsquare}

\emph{Status: proved (definition + validation), not yet instantiated per
candidate at the level of an explicit exhibited $L_2$.}

We formalize monic order-2 and order-3 operators as structures
\lstinline{DiffOp2}, \lstinline{DiffOp3} carrying \lstinline{Polynomial ℚ}
coefficients, and the symmetric-square map
\lstinline{symSquare : DiffOp2 → DiffOp3} implementing
\eqref{eq:symsquare-formula} verbatim. The relation
\lstinline{IsSymSquareOf L3 L2 := (L3 = symSquare L2)} is a \emph{concrete}
operator equality --- every coefficient of $L_3$ is pinned as an explicit
polynomial function of $L_2$'s coefficients --- and is therefore falsifiable
per candidate, in contrast to an existence statement.

Formula~\eqref{eq:symsquare-formula} is validated, independently of the
literature, against two examples worked from first principles and proved as
golden tests:
\begin{itemize}
  \item $L_2 = D^2+1$ (harmonic oscillator; solutions $\sin t,\cos t$) must
    give $\mathrm{Sym}^2(L_2) = D^3+4D$, since $\{\sin^2,\sin\cos,\cos^2\}$
    span the same space as $\{1,\cos 2t,\sin 2t\}$, annihilated by
    $D(D^2+4)$;
  \item $L_2 = D^2+D$ (solutions $1,e^{-t}$) must give
    $\mathrm{Sym}^2(L_2)=D^3+3D^2+2D$, since $\{1,e^{-t},e^{-2t}\}$ are
    annihilated by $D(D+1)(D+2)$.
\end{itemize}
Both reduce, after unfolding the definition, to a single \texttt{simp} call.

% ════════════════════════════════════════════════════════════════════
\section{Main result: the generic self-adjointness theorem}
\label{sec:w0}

\emph{Status: proved. This is the paper's principal new contribution.}

\subsection{Statement}

Converting~\eqref{eq:cooper-theta} to standard ($D=d/dz$) form via the
substitutions $\theta = zD$, $\theta^2 = z^2D^2+zD$,
$\theta^3=z^3D^3+3z^2D^2+zD$ and collecting powers of $D$ gives
$L = p_3 D^3 + p_2 D^2 + p_1 D + p_0$ with
\begin{align}
  p_3 &= z^3\bigl(1-2az+cz^2\bigr), &
  p_2 &= 3z^2\bigl(1-3az+2cz^2\bigr), \\
  p_1 &= z\bigl(1-(6a+2b)z+(7c+d)z^2\bigr), &
  p_0 &= z\bigl(-b+(c+d)z\bigr).
\end{align}
These coefficients were cross-checked three independent ways before being
trusted: (a) by hand, substituting the $\theta\to D$ identities above into
the \emph{already machine-proved} collected $\theta$-form
(\lstinline{cooperThetaOperator_eq}, Section~\ref{sec:theta}); (b) against an
independent symbolic re-derivation (computer algebra); (c) against a second,
fully independent symbolic re-derivation performed by a second reviewer
without access to the first derivation (the ``two-model'' verification
protocol used throughout this project for any claim not directly
kernel-checked). All three agree exactly.

The monic coefficients $a_i = p_i/p_3$ would be needed to state the
Almkvist--van Straten criterion~\eqref{eq:W-criterion} directly, but $p_3$ is
a non-unit polynomial, so this route re-introduces the rational-function
obstruction of Section~\ref{sec:theta}. Instead, since $p_3\neq 0$ for every
parameter choice (its constant term is~1), the criterion $W\equiv 0$ is
equivalent, in the fraction field $\mathbb{Q}(a,b,c,d,z)$, to the vanishing of
its numerator after clearing denominators by $27p_3^3$:
\begin{align}
  \label{eq:Pcleared}
  P_{\mathrm{cleared}} \;:=\;\;
    & 9\bigl(p_2'' p_3^2 - p_2 p_3 p_3'' - 2p_2' p_3 p_3' + 2p_2 (p_3')^2\bigr) \nonumber\\
  {}+{} & 18\, p_2\bigl(p_2' p_3 - p_2 p_3'\bigr) \;+\; 4p_2^3 \;+\; 54\, p_0 p_3^2 \nonumber\\
  {}-{} & 18\, p_1 p_2 p_3 \;-\; 27\, p_3\bigl(p_1' p_3 - p_1 p_3'\bigr),
\end{align}
a genuine element of $\mathbb{Q}[a,b,c,d][z]$ --- no division required. We
prove:

\begin{theorem}[Generic self-adjointness, \lstinline{P_cleared_eq_zero}]
\label{thm:w0}
For \emph{every} $a,b,c,d\in\mathbb{Q}$, $P_{\mathrm{cleared}}(a,b,c,d,z) = 0$
as an element of $\mathbb{Q}[z]$.
\end{theorem}

By the equivalence above, Theorem~\ref{thm:w0} states that the Cooper
$\theta$-operator family satisfies the Almkvist--van Straten $W\equiv 0$
criterion \emph{identically}, for every choice of the four parameters, and in
particular for $s_7$, $s_{10}$, and $s_{18}$ simultaneously ---
Corollaries~\ref{cor:s7}--\ref{cor:s18} below are immediate substitutions,
requiring no further proof.

\begin{corollary}\label{cor:s7} $P_{\mathrm{cleared}}(13,4,-27,3,z)=0$.\end{corollary}
\begin{corollary}\label{cor:s10} $P_{\mathrm{cleared}}(6,2,-64,4,z)=0$.\end{corollary}
\begin{corollary}\label{cor:s18} $P_{\mathrm{cleared}}(14,6,192,-12,z)=0$.\end{corollary}

\begin{remark}
Theorem~\ref{thm:w0} establishes that $L_3$ \emph{admits a} symmetric-square
factorization (existence, via the classical equivalence
$W\equiv0 \Leftrightarrow L_3=\mathrm{Sym}^2(L_2)$ for some $L_2$); it is a
structural statement about the entire Cooper ansatz, and, following an
observation made during this project's review process, does \emph{not by
itself} distinguish among the three candidates $s_7,s_{10},s_{18}$, since it
holds for all of them (and indeed for every choice of $a,b,c,d$) uniformly.
Exhibiting the specific order-2 operator $L_2$ for a given candidate, and
verifying \lstinline{IsSymSquareOf} (Section~\ref{sec:symsquare}) against it
directly, remains separate, not-yet-formalized work.
\end{remark}

\subsection{Proof discipline}

In keeping with the ``kernel is the judge'' principle underlying this
development, the derivatives in~\eqref{eq:Pcleared} are \emph{not}
hand-computed closed forms substituted into the Lean statement; they are
literal applications of mathlib's \lstinline{Polynomial.derivative} to the
polynomials $p_0,\dots,p_3$ as defined in Lean, so that the kernel itself
performs and checks every differentiation step. An opaque, hand-transcribed
``derivative'' definition would let an arithmetic slip in the transcription
pass \texttt{ring} undetected; using the genuine formal derivative closes
that gap. The proof of Theorem~\ref{thm:w0} is a single \texttt{simp}
normalization (unfolding \lstinline{Polynomial.derivative} across sums,
products, and powers, and normalizing scalar multiples of \(X^n\) introduced
by the power rule) followed by \texttt{ring}, which decides the resulting
polynomial identity. The complete proof is reproduced in
Appendix~\ref{app:symsquarec3b}.

% ════════════════════════════════════════════════════════════════════
\section{Growth-rate correction}
\label{sec:growth}

\emph{Status: proved (correction of a previously false claim).}

An earlier version of this repository recorded, as open conjectures, the
statements that $s_7$ and $s_{10}$ are bounded by a polynomial in $n$. Direct
numerical inspection of the exact-arithmetic ratio $u(n{+}1)/u(n)$ shows it
increasing without bound over the first 25 terms (past 25 for $s_7$, past 15
for $s_{10}$), the signature of exponential rather than polynomial growth ---
consistent with the standard asymptotic theory of D-finite (holonomic)
sequences, whose growth rate is set by the recurrence's nearest finite
singularity, generically giving $\rho^n n^\alpha$ growth rather than a
polynomial bound. We record the corrected, proved statements:

\begin{theorem}[\lstinline{s7_not_polynomially_bounded}, \lstinline{s10_not_polynomially_bounded}]
Neither $s_7$ nor $s_{10}$ is bounded by any polynomial: there is no pair
$(c,k)\in\mathbb{N}^2$ with $s_7(n)\le c(n+1)^k$ (resp.\ $s_{10}$) for all $n$.
\end{theorem}

The proof exhibits, for each sequence, a clean single-term lower bound via the
central binomial coefficient ($s_{10}(2n)\ge \binom{2n}{n}^4$, taking the
$j=n$ term of the defining sum; $s_7(n) \ge \binom{2n}{n}^2$, taking the $k=n$
term) and combines mathlib's exponential lower bound on the central binomial
coefficient with the standard fact that any fixed-degree polynomial is
eventually dominated by any exponential of base $>1$. As documented in the
source, this is a \emph{correction}, not a weakening: the original claim was
false as stated, and the false open goals were removed from the repository in
the same change that introduced this disproof.

% ════════════════════════════════════════════════════════════════════
\section{Integrality and the scope of recurrence verification}
\label{sec:integrality}

\emph{Status: mixed --- see below for exactly which claims are proved versus
open.}

Both $s_7(n)$ and $s_{10}(n)$ are integers for every $n$ \emph{by
construction} (they are defined as \lstinline{Finset.sum}s of natural-number
values), so integrality is definitional rather than a separate theorem.

The substantive open question is whether the closed-form binomial sums of
Table~\ref{tab:params} satisfy the three-term recurrence
\eqref{eq:cooper-recurrence} for \emph{every} $n$. This is mathematically
established in the literature (via a Wilf--Zeilberger / creative-telescoping
certificate, per Gorodetsky~\cite{Gorodetsky2023}) but is \textbf{not proved
in this repository}. We record the precise status:

\begin{itemize}
  \item \emph{Proved:} the recurrence holds at $n=1$, by direct kernel
    computation (\lstinline{s7_recurrence_at_one}, \lstinline{s10_recurrence_at_one}).
  \item \emph{Checked-but-unproved:} the recurrence holds for
    $n=1,\dots,18$, verified by \texttt{native\_decide}
    (Section~\ref{sec:sequences}'s disclosure applies). This is evidence for,
    not a proof of, the universal statement.
  \item \emph{Open:} the $\forall n$ statement itself
    (\lstinline{open_goal_recurrence_s7}, \lstinline{open_goal_recurrence_s10}
    in the repository's machine-readable open-goal register). Three genuinely
    different proof strategies were attempted and are documented as failed in
    the source (induction on $n$ with the defining sum's shifting index range;
    kernel \texttt{decide} at the $\forall n$ level; a search for a matching
    mathlib lemma) before the statement was correctly classified as open
    rather than continuing unbounded proof search, per this project's
    ``three-strikes'' engineering discipline. Closing it requires formalizing
    the underlying Wilf--Zeilberger certificate as an auxiliary lemma, which
    has not yet been transcribed into this repository.
\end{itemize}

We flag this explicitly because it would be easy, and dishonest, to elide the
distinction between ``verified for $n\le 18$'' and ``proved for all $n$.''

% ════════════════════════════════════════════════════════════════════
\section{Scope boundary: what this report does not claim}
\label{sec:scope}

\emph{This section exists to make an explicit negative statement, as a matter
of intellectual honesty.}

The host repository contains, in addition to the material above, exploratory
Lean code exploring a purely speculative physical framework (an F-theory-style
moduli-stabilization argument and a chameleon-screening argument concerning a
hypothetical ultralight axion near the supermassive black hole M87*). None of
that material is used by, or required for, any result in this paper, and none
of it is reported here, for two concrete reasons documented in the
repository's own source:
\begin{enumerate}
  \item one axiom underlying part of that material is explicitly
    self-documented in its own docstring as \emph{vacuous} (it asserts the
    existence of a real number satisfying a bound, witnessed trivially by
    $1$, and is stated to encode no actual pipeline data pending a
    certified data artifact that does not yet exist);
  \item the numerical parameters entering the remaining arguments (e.g.\ an
    assumed density ratio and bare coupling for the black hole M87*) are
    stipulated constants chosen to produce a target inequality, not values
    derived from, or fit to, any observational data or first-principles
    physical model.
\end{enumerate}
A separate internal memorandum, not styled or intended as a scientific report,
catalogs these gaps for the repository maintainer's review. We mention their
existence here solely so that a reader of this paper's source repository is
not misled by documentation elsewhere in that repository into believing this
paper's mathematical results establish, or depend on, any physical claim.

% ════════════════════════════════════════════════════════════════════
\section{Reproducibility}
\label{sec:repro}

\begin{itemize}
  \item \textbf{Language/tooling:} Lean~4 (toolchain \texttt{leanprover/lean4:v4.32.0}),
    mathlib4 pinned at commit \texttt{3dffaf2f18b47d11948f6390838ea6f2ae662aaf}~\cite{mathlib}.
  \item \textbf{Build:} \texttt{lake build Agora} builds the full development
    (3106 jobs at the time of writing); \texttt{lake build Tests} runs the
    golden numeric tests. Both complete with zero errors and zero warnings on
    the files discussed in this paper.
  \item \textbf{Sorry/axiom audit:} zero \lstinline{sorry} and zero axioms
    outside the quarantined \texttt{Axioms/} directory in any file cited in
    this paper; the one repository-wide axiom (Section~\ref{sec:scope}) does
    not occur in, and is not reachable from, any result reported here.
  \item \textbf{Open-goal export:} \texttt{python3 scripts/export\_open\_goals.py}
    regenerates a machine-readable list of open goals; at the time of writing
    it lists exactly the two items of Section~\ref{sec:integrality} and no
    others.
  \item \textbf{Source:} the files discussed in this paper are reproduced
    verbatim in Appendix~\ref{app:source}; the canonical, buildable source
    lives in the project git repository at commit
    \texttt{60c1374} on branch \texttt{main}.
\end{itemize}

% ════════════════════════════════════════════════════════════════════
\section{Discussion and outlook}

The results above formalize a modest but precisely delimited slice of the
theory of sporadic Ap\'ery-like sequences: sourced encodings, operator-order
facts, a validated symmetric-square API, and --- the paper's central
contribution --- a single generic proof that the entire Cooper family
satisfies the Almkvist--van Straten self-adjointness criterion, sidestepping
a genuine formalization obstruction (the absence of a \lstinline{RatFunc}
differentiation API at the pinned mathlib commit) by working with a
cleared-denominator polynomial identity instead of the naive monic
formulation. We regard the explicit, disclosed boundary of what remains open
(Section~\ref{sec:integrality}) as no less a contribution than the proved
results themselves: it converts an implicit assumption into an auditable,
named proof obligation for future work, specifically the formalization of a
Wilf--Zeilberger certificate for the two remaining $\forall n$ recurrence
statements.

We deliberately do not speculate here about the physical relevance of these
structures (Section~\ref{sec:scope}); any such discussion belongs in a
separate report once the corresponding claims meet the same standard of
verification applied throughout this paper.

% ════════════════════════════════════════════════════════════════════
\bibliographystyle{plain}
\bibliography{references}

% ════════════════════════════════════════════════════════════════════
\appendix
\section{Complete Lean 4 source}
\label{app:source}

The following listings reproduce, verbatim and in full, every Lean file whose
results are reported in this paper. They are included by direct file
inclusion (\texttt{lstinputlisting}) from the accompanying repository, so
there is no risk of transcription drift between this document and the
buildable source; the line numbers shown are the files' own.

\subsection{\texttt{Agora/Sequences/CooperRecurrences.lean}}
\label{app:cooper-recurrences}
\lstinputlisting[language=Lean4]{../Agora/Sequences/CooperRecurrences.lean}

\subsection{\texttt{Agora/Sequences/ThetaOperators.lean}}
\label{app:theta-operators}
\lstinputlisting[language=Lean4]{../Agora/Sequences/ThetaOperators.lean}

\subsection{\texttt{Agora/SymSquare.lean}}
\label{app:symsquare}
\lstinputlisting[language=Lean4]{../Agora/SymSquare.lean}

\subsection{\texttt{Agora/Swampland/SymSquareC3b.lean}}
\label{app:symsquarec3b}
\lstinputlisting[language=Lean4]{../Agora/Swampland/SymSquareC3b.lean}

\subsection{\texttt{Agora/Sequences/GrowthBounds.lean}}
\label{app:growth-bounds}
\lstinputlisting[language=Lean4]{../Agora/Sequences/GrowthBounds.lean}

\subsection{\texttt{Agora/Sequences/Integrality.lean}}
\label{app:integrality}
\lstinputlisting[language=Lean4]{../Agora/Sequences/Integrality.lean}

\subsection{\texttt{Tests/CooperSequences.lean} (golden numeric tests)}
\label{app:tests}
\lstinputlisting[language=Lean4]{../Tests/CooperSequences.lean}

\end{document}
